\begin{prob}
    In each of the problems below compute the average case time complexity (or
    expected time) of the given code. State your answer using asymptotic
    notation in the simplest terms possible. Show your work for this problem by stating what the different
    cases are, the probability of each case, and how long each case takes. Also
    show the calculation of the expected time.

    \begin{subprobset}
        \begin{subprob}
            \begin{minted}[autogobble]{python}
                def foo(n):
                    # randomly choose a number between 0 and n-1 in constant time
                    k = np.random.randint(n)

                    if k < np.sqrt(n):
                        for i in range(n**2):
                            print(i)
                    else:
                        print('Never mind...')

            \end{minted}

            \begin{soln}
                $\Theta(n\sqrt{n})$.

                We can think of the two cases here as being 1) $k$ is randomly chosen
                to be less than $\sqrt{n}$, and 2) $k$ is randomly chosen to be greater.
                The probability of the first case is $\sqrt{n} / n = 1/\sqrt{n}$, and
                the probability of the second case is $(n - \sqrt{n})/n = 1 - 1/\sqrt{n}$.

                In the first case, quadratic time (that is, $\Theta(n^2)$ time)
                is taken since we must loop over \python{range(n**2)}. In the
                second case, constant time is taken since we print
                \python{"Never mind..."} and exit.

                Therefore, the average time taken is:

                \begin{align*}
                    \frac{1}{\sqrt{n}} \Theta(n^2) + \left(1 - \frac{1}{\sqrt{n}}\right) \Theta(1)
                    =
                        \Theta(n\sqrt{n}) + \Theta(1)
                    =
                        \Theta(n\sqrt{n})
                \end{align*}

            \end{soln}
        \end{subprob}


        \begin{subprob}
            \begin{minted}[autogobble]{python}
                def bogosearch(numbers, target):
                    """search by randomly guessing. `numbers` is an array of n numbers"""
                    n = len(numbers)

                    while True:
                        # randomly choose a number between 0 and n-1 in constant time
                        guess = np.random.randint(n)
                        if numbers[guess] == target:
                            return guess
            \end{minted}

            In this part, you may assume that the numbers are distinct and that the
            target is in the array.

            Hint: if $0 < b < 1$, then $\sum_{p = 1}^\infty p \cdot b^{p-1} = \frac{1}{(1 - b)^2}$.

            \begin{soln}
                $\Theta(n)$.

                Maybe the best way to divide this into cases is to consider Case 1 to
                be when we guess the location of the target on the very first try;
                Case 2 is when we guess the location of the target incorrectly on the
                first try, but get it right on the second; Case 3 is when we get the
                first two guesses wrong but get the third guess correct, and so on.
                In general, Case $\alpha$ is when we made incorrect guesses on the first
                $\alpha - 1$ iterations, but a correct guess on the $\alpha$th iteration,
                therefore returning.

                The probability of Case $\alpha$ is the probability of guessing wrong $\alpha - 1$
                times followed by guessing correctly. Since these guesses are all statistically
                independent, and since the probability of any one guess being wrong is
                $(n - 1)/n$ and the probability of it being right is $1/n$, we find that
                the probability of Case $\alpha$ is:
                \[
                    \underbrace{
                        \left( \frac{n-1}{n} \right)
                        \left( \frac{n-1}{n} \right)
                        \cdots
                        \left( \frac{n-1}{n} \right)
                    }_{\alpha - 1 \text{ times}}
                    \cdot \frac{1}{n}
                    =
                    \left(
                        \frac{n-1}{n}
                    \right)^{\alpha-1} \frac{1}{n}
                    =
                    \left(
                        1 - \frac{1}{n}
                    \right)^{\alpha-1} \frac{1}{n}
                \]

                The time taken on any iteration is constant; let's call it $c$. In case $\alpha$
                we make $\alpha$ iterations, and so we take $c \alpha$ time in total.

                There are actually infinitely many cases here, since it is possible that
                we are consistently unlucky and never guess the right entry. Therefore,
                our summation will actually be an infinite series.

                Therefore the expected time over all possible cases is:
                \begin{align*}
                    \sum_{\alpha = 1}^\infty
                        P(\text{case } \alpha)
                        \cdot
                        T(\text{case } \alpha)
                        =
                    \sum_{\alpha = 1}^\infty
                    \left(
                        1 - \frac{1}{n}
                    \right)^{\alpha-1} \frac{1}{n} \cdot c \alpha
                    &=
                        \frac{c}{n}
                        \sum_{\alpha = 1}^\infty
                        \alpha
                        \left(
                            1 - \frac{1}{n}
                        \right)^{\alpha-1}
                        \\
                    \intertext{We can now use the hint. Letting $b = 1 - 1/n$, we have:}
                    &=
                        \frac{c}{n}
                        \sum_{\alpha = 1}^\infty
                        \alpha
                        b^{\alpha - 1}
                        \\
                    &=
                        \frac{c}{n} \cdot \frac{1}{(1 - b)^2}
                    \intertext{Now substituting back in for $b = 1 - 1/n$:}
                    &=
                    \frac{c}{n} \cdot \frac{1}{\left[1 - (1 - 1/n)\right]^2}
                        \\
                    &=
                        \frac{c}{n} \cdot \frac{1}{1/n^2}
                        \\
                    &=
                        \frac{c}{n} \cdot n^2
                        \\
                    &=
                        cn
                \end{align*}

                So the average case time complexity is $\Theta(n)$.

            \end{soln}


        \end{subprob}
    \end{subprobset}
\end{prob}
